% STATUS: verified
\chapter{Glossary and abbreviations}\label{app:glossary}

The words this document uses for the system's own things, the labels the console prints for
each status in its two languages, and the abbreviations of acts and institutions that a reader
from outside Lithuania may not know.

\section{The system's vocabulary}\label{sec:glossary:terms}

\begin{xltabular}{\textwidth}{@{}L{36mm}Y@{}}
\toprule
\hd{Term} & \hd{Meaning in this system} \\
\midrule
\endfirsthead
\toprule
\hd{Term} & \hd{Meaning in this system} \\
\midrule
\endhead
\bottomrule
\endlastfoot
Act & A legal text of the corpus: a regulation, a directive, a law, a resolution, an order, a
  guideline, a licensing page. A standard is kept the same way and called a standard. \\
Action plan & The list of actions a report proposes; every open finding is under exactly one
  of them. \\
Action register & A company's list of actions that lives across reports. Code moves it only
  when a person signs a report off or the signature is taken back; between those, a person
  keeps each action up by hand. \\
Applicability & The verdict on whether an act binds a company: applies, does not apply, or
  uncertain (the console says \enquote{to be confirmed}). \\
Assessment (of an act) & What a report does with an act: \emph{assessed}; \emph{upcoming} ---
  assessed as preparation because it binds from a later date; \emph{unconfirmed} --- not
  assessed until the company answers the open questions; \emph{indirect} --- a directive,
  which reaches the company through the national law that transposes it, or an act that only
  organises an authority, which reaches it through the authority's decisions. \\
Catalogue & The list of acts the product ships with, and their sources. \\
Chunk & A passage of an act or of a company document, stored with its embedding for search. \\
Citation check & The check of an article a document cites: does the provision exist, and does
  it regulate what the document says it does. Inside a report it is the \emph{citation audit}. \\
Consequence & What a breach of a finding's provision can cost, joined by code from the act's
  sanction tiers: the ceiling the act provides, not a forecast. It is shown only for a gap. \\
Corpus (legal corpus) & All the acts with their text, in versions, parsed into provisions. \\
Embedding & A vector that represents a passage of text, so that passages of similar meaning
  can be found. \\
Finding & The result for one obligation in one report: a status, an assessment, the evidence
  quoted, a recommendation, the consequence, and --- if a reviewer wrote one --- a cause. \\
Gateway & The faculty's model gateway: in the deployment described, the only way the system
  reaches a language model. \\
Inventory & What one company document establishes, read once into a structured record:
  controls with quotations, roles, registers, approval, the provisions it cites. \\
Job, step & A unit of background work and its parts. A report is made by the job
  \code{report.analyze} in eight steps, nine when the company's own requirements are assessed
  too; a report started with the scope review runs as two such jobs, one before the approval
  and one after it. \\
Judge & A second model call that checks the output of a first against the source: a separate
  call with a prompt of its own, not necessarily another model (on the live system every judge runs on
  \model{gpt-6.1-sol}). It doubts or supports; it approves nothing. \\
Library & All the documents of one company. \\
Obligation & A duty an act imposes, written down once with who owes it, under what condition,
  how gravely, and the provisions it comes from. \\
Own requirement & A requirement a company set itself in one of its policies, read out so that
  a report can check whether the procedures carry it out. \\
Path card & In this document: the table of ten rows that describes one place where a model is
  called. \\
Prompt, prompt version & The stored instructions for one kind of model call, and one immutable
  edition of them. \\
Provision, unit & An article, paragraph, point or clause of an act; as a stored row, a unit. \\
Register of obligations & All the obligations of all the acts, with their review status. \\
Report & The compliance report of one company at one time. \\
Scope review & The stop in the middle of a report at which a person approves the list of acts
  before the costly part runs. \\
Severity & How much an obligation matters: critical, high, medium or low. It belongs to the
  obligation, not to a finding. \\
Sign-off & A person's confirmation of a finished report. It moves the action register. \\
Tier (sanction tier) & One rule of an act's sanction provisions: which infringements, what
  ceiling of fine, what other measures. \\
Version & One stored edition of an act. \\
Work product & A document of a compliance programme drafted by the system: the policy, the
  structure with the process and the responsibility matrix, the risk register. \\
Worker & The process that runs jobs. \\
\end{xltabular}

\section{Statuses and their labels}\label{sec:glossary:labels}

The code stores a status under an English key; the console prints a label in the language of
the interface. This document uses the stored names, with \enquote{not evidenced} for
\code{unknown}.

\begin{xltabular}{\textwidth}{@{}L{26mm}L{26mm}L{40mm}Y@{}}
\toprule
\hd{Of} & \hd{Stored as} & \hd{Lithuanian label} & \hd{English label} \\
\midrule
\endfirsthead
\toprule
\hd{Of} & \hd{Stored as} & \hd{Lithuanian label} & \hd{English label} \\
\midrule
\endhead
\bottomrule
\endlastfoot
A finding & \code{met} & \lt{Atitinka} & Met \\
 & \code{partial} & \lt{Iš dalies} & Partial \\
 & \code{missing} & \lt{Neatitinka} & Not met \\
 & \code{unknown} & \lt{Neįrodyta} & Not evidenced \\
 & \code{not_applicable} & \lt{Netaikoma} & Not applicable \\
\midrule
Severity & \code{critical} & \lt{Kritinė} & Critical \\
 & \code{high} & \lt{Aukšta} & High \\
 & \code{medium} & \lt{Vidutinė} & Medium \\
 & \code{low} & \lt{Žema} & Low \\
\midrule
Applicability & \code{applies} & \lt{Taikomas} & Applies \\
 & \code{not_applies} & \lt{Netaikomas} & Does not apply \\
 & \code{uncertain} & \lt{Tikslintinas} & To be confirmed \\
\midrule
A report & \code{queued} & \lt{Laukia eilėje} & Queued \\
 & \code{running} & \lt{Vykdoma} & Running \\
 & \code{scope_review} & \lt{Laukia apimties patvirtinimo} & Awaiting scope approval \\
 & \code{done} & \lt{Baigta} & Finished \\
 & \code{failed} & \lt{Nepavyko} & Failed \\
 & \code{cancelled} & \lt{Atšaukta} & Cancelled \\
\midrule
An obligation & \code{proposed} & \lt{Pasiūlyta} & Proposed \\
 & \code{approved} & \lt{Patvirtinta} & Approved \\
 & \code{rejected} & \lt{Atmesta} & Rejected \\
\midrule
The judge's verdict & \code{supported} & \lt{Pagrįsta} & Supported \\
 & \code{partly} & \lt{Iš dalies} & Partly \\
 & \code{unsupported} & \lt{Nepagrįsta} & Unsupported \\
\midrule
An action & \code{open} & \lt{Atviras} & Open \\
 & \code{in_progress} & \lt{Vykdomas} & In progress \\
 & (derived) & \lt{Vėluojantis} & Late \\
 & \code{done} & \lt{Įvykdytas, laukia patikros} & Done, awaiting re-test \\
 & \code{closed} & \lt{Uždarytas} & Closed \\
\midrule
A work product & \code{queued} & \lt{Laukia eilėje} & Queued \\
 & \code{running} & \lt{Generuojama} & Generating \\
 & \code{draft} & \lt{Projektas} & Draft \\
 & \code{approved} & \lt{Patvirtinta} & Approved \\
 & \code{failed} & \lt{Nepavyko} & Failed \\
 & \code{cancelled} & \lt{Atšaukta} & Cancelled \\
\end{xltabular}

\section{Abbreviations of acts}\label{sec:glossary:acts}

The system names an act by its short title, which for a Lithuanian law is the abbreviation
Lithuanian lawyers use. \Cref{app:acts} lists every act with its subject.

\begin{xltabular}{\textwidth}{@{}L{20mm}YY@{}}
\toprule
\hd{Short title} & \hd{In Lithuanian} & \hd{In English} \\
\midrule
\endfirsthead
\toprule
\hd{Short title} & \hd{In Lithuanian} & \hd{In English} \\
\midrule
\endhead
\bottomrule
\endlastfoot
ADTAĮ & \lt{Asmens duomenų teisinės apsaugos įstatymas} & Law on the Legal Protection of Personal Data \\
BĮ & \lt{Bankų įstatymas} & Law on Banks \\
EPEPĮĮ & \lt{Elektroninių pinigų ir elektroninių pinigų įstaigų įstatymas} & Law on Electronic Money and Electronic Money Institutions \\
FĮĮ & \lt{Finansų įstaigų įstatymas} & Law on Financial Institutions \\
KRĮ & \lt{Kriptoturto rinkų įstatymas} & Law on Markets in Crypto-Assets \\
KSĮ & \lt{Kibernetinio saugumo įstatymas} & Law on Cyber Security \\
KUĮ & \lt{Kredito unijų įstatymas} & Law on Credit Unions \\
LBĮ & \lt{Lietuvos banko įstatymas} & Law on the Bank of Lithuania \\
MĮ & \lt{Mokėjimų įstatymas} & Law on Payments \\
MĮĮ & \lt{Mokėjimo įstaigų įstatymas} & Law on Payment Institutions \\
NTSKĮ & \lt{Su nekilnojamuoju turtu susijusio kredito įstatymas} & Law on Credit Related to Immovable Property \\
PAĮ & \lt{Pranešėjų apsaugos įstatymas} & Law on the Protection of Whistleblowers \\
PPTFPĮ & \lt{Pinigų plovimo ir teroristų finansavimo prevencijos įstatymas} & Law on the Prevention of Money Laundering and Terrorist Financing \\
VKĮ & \lt{Vartojimo kredito įstatymas} & Law on Consumer Credit \\
\midrule
AMLAR & & Regulation (EU) 2024/1620 establishing the Authority for Anti-Money Laundering and Countering the Financing of Terrorism \\
AMLD6 & & Directive (EU) 2024/1640 on the mechanisms Member States put in place against money laundering and terrorist financing \\
AMLR & & Regulation (EU) 2024/1624 on the prevention of money laundering and terrorist financing \\
CCD2 & & Directive (EU) 2023/2225 on credit agreements for consumers \\
CRA & & Regulation (EU) 2024/2847, the Cyber Resilience Act \\
DORA & & Regulation (EU) 2022/2554 on digital operational resilience for the financial sector \\
DSA & & Regulation (EU) 2022/2065, the Digital Services Act \\
ECSPR & & Regulation (EU) 2020/1503 on European crowdfunding service providers \\
eIDAS~2 & & Regulation (EU) No 910/2014 on electronic identification and trust services, as amended in 2024 \\
EMD2 & & Directive 2009/110/EC on electronic money institutions \\
GDPR & & Regulation (EU) 2016/679, the General Data Protection Regulation \\
MAR & & Regulation (EU) No 596/2014 on market abuse \\
MiCA & & Regulation (EU) 2023/1114 on markets in crypto-assets \\
NIS2 & & Directive (EU) 2022/2555 on a high common level of cybersecurity across the Union \\
PSD2 & & Directive (EU) 2015/2366 on payment services in the internal market \\
TFR & & Regulation (EU) 2023/1113 on information accompanying transfers of funds and certain crypto-assets \\
\end{xltabular}

\section{Other abbreviations}\label{sec:glossary:other}

\begin{xltabular}{\textwidth}{@{}L{24mm}Y@{}}
\toprule
\hd{Abbreviation} & \hd{Stands for} \\
\midrule
\endfirsthead
\toprule
\hd{Abbreviation} & \hd{Stands for} \\
\midrule
\endhead
\bottomrule
\endlastfoot
CASP & Crypto-asset service provider. \\
CELEX & The identifier of a legal document of the European Union (for example
  \code{32023R1114} for MiCA). \\
CELLAR & The repository of the Publications Office of the European Union, from which the
  system fetches EU acts. \\
EBA & European Banking Authority. \\
e-tar & \lt{Teisės aktų registras}, Lithuania's register of legal acts, from which the system
  fetches the Lithuanian laws, the Bank of Lithuania's resolutions and the FNTT's order; the
  Bank of Lithuania's guidelines and licensing pages come from its own website. \\
FNTT & \lt{Finansinių nusikaltimų tyrimo tarnyba}, the Financial Crime Investigation Service. \\
HNSW & Hierarchical navigable small world: the kind of index used for nearest-neighbour search
  over embeddings. \\
ISMS & Information security management system. \\
ITS, RTS & Implementing and regulatory technical standards: Commission regulations that
  specify how a provision of a parent regulation is carried out (\enquote{level~2}). \\
LB & \lt{Lietuvos bankas}, the Bank of Lithuania, the supervisor of the financial market. \\
MCP & Model Context Protocol: a protocol through which the chat assistant can use tools
  offered by an external server an administrator has registered. \\
ML/TF & Money laundering and terrorist financing. \\
RACI & Responsible, accountable, consulted, informed: the four ways a role takes part in a
  step of a process. \\
RRF & Reciprocal rank fusion: the way the results of vector search and of full-text search are
  merged into one ranking. \\
SNI & Server Name Indication: the host name a client names when it opens an encrypted
  connection. The system's only way out to the internet allows connections by this name. \\
\end{xltabular}
